\chapter{Riemannian Geometry}
\label{chap:riemannian-geometry}

% Closed question: What reference structure lets distinct observers compare
% their counts as one physical record?

\begin{chapteressay}

Two observers measure the same physical quantity. One reports the
length in meters; the other reports it in feet. To compare their
records, a conversion rate must be named: 1 meter equals about 3.281
feet. The conversion is the chapter's reference structure: the
shared piece of information that makes the two records commensurable.

Without the conversion, the two records are not directly comparable.
The first observer's ``5 meters'' and the second observer's ``18
feet'' could refer to roughly the same physical length, or to two
quite different lengths; without the conversion, the comparison is
ambiguous. With the conversion, the two records can be brought to a
common scale, and the comparison becomes admissible.

The chapter is about the reference structure that lets distinct
observers compare their counts as one physical record. The
reference is the calibration certificate; the calibration governs
the conversion; the conversion makes the comparison admissible.

The chapter develops the reference structure through five
mathematical objects. The first is metric from counting: a metric
can be constructed from counting alone, without any geometric
assumption. The Hamming distance is the canonical example: count
the positions at which two strings differ, and that count is the
distance.

The second is the Cartesian construction: Descartes's product of
perpendicular rulers gives the plane, with the Pythagorean formula
forced by the invariance of distance under rotation. The
construction is purely algebraic: a product of one-dimensional
metric records gives a two-dimensional metric record.

The third is the calibration certificate: a named record of the
reference conditions under which two observers' counts can be
compared. The \texttt{EKG} typeclass in Lean is the canonical
algebraic version: a private reference scale (the
\texttt{Reading}) accessed only through a public comparison
interface (the \texttt{boolean} method).

The fourth is the metric tensor: the local reference structure on
a curved manifold. The tensor varies with position; the
variation captures the manifold's curvature; the local distance
is determined by the local tensor.

The fifth is truth as carried structure: a theorem's truth is a
carried certificate (the proof); removing the proof removes the
certificate; the truth is portable as the proof is portable.

The chapter's ground example is the ratio of two counts in
different units. The conversion rate is the calibration; the
conversion translates one count into another; the translated
counts are admissibly comparable.

The chapter's second concrete example is the currency exchange
desk. The exchange rate between two currencies is the chapter's
reference structure. The rate translates dollars into euros (or
yen, or pounds); the translation makes two travelers' wealths
comparable in a common currency. The consistency conditions on the
rate matrix (identity, inverse, triangle) are the calibration
requirements. Arbitrage enforces consistency.

The chapter's closed question is:

\begin{center}
\emph{What reference structure lets distinct observers compare
their counts as one physical record?}
\end{center}

The chapter's answer is: the calibration certificate plus the
metric. The calibration is the named reference; the metric is
the structural rule for distance computation; the two together
let observers compare counts.

Subsequent chapters refine this. Chapter 9 introduces curvature:
the geometric quantity that records the failure of transport to
commute. Curvature is computed from the metric tensor's
derivatives. Chapter 10 introduces invariance: what survives
admissible relabelings of the coordinate system.

\end{chapteressay}

\section{Metric from Counting}
\label{se:metric-from-counting}

Two DNA strands are compared. The first is ACGT; the second is ACCT.
The Hamming distance between them is one: they differ at exactly one
position (the third). The Hamming distance is a metric on the space
of strings of equal length: it satisfies the triangle inequality, it
is symmetric, it is positive definite.

The Hamming distance is constructed entirely from counting. The
operation: align the two strings position by position; count the
positions at which they differ; the count is the distance. No
geometry is involved; no number line is required. The distance
emerges from the position-by-position comparison and the count of
differences.

This is the chapter's first claim. A metric can be constructed
from counting alone. The counting is the comparison operation
applied to corresponding components of the two records; the
count is the metric. No reference to a continuous quantity is
needed.

The Hamming distance is the chapter's simplest example. More
elaborate counting-based metrics include:
\begin{itemize}
\item Levenshtein (edit) distance: the minimum number of
insertions, deletions, or substitutions transforming one string
into the other.
\item Jaccard distance: $1 - |A \cap B| / |A \cup B|$ for sets $A$
and $B$. The distance counts the symmetric difference relative
to the union.
\item Cosine distance: $1 - \cos(\theta)$ for the angle $\theta$
between two vectors. Despite the appearance of geometry, the
cosine is computed from inner products of counts.
\end{itemize}

Each metric is built from counting. The geometric interpretation
(if any) is downstream of the counting; the metric exists in the
combinatorial structure of the records.

The chapter's claim is that this construction extends to physical
measurement. Every measurement is fundamentally a count: the
yardstick records how many ticks fit; the clock records how many
tocks have passed; the thermometer records the position of a
column relative to graduated marks. The metric on physical
quantities is built from these counts.

\section{Descartes: The Product of Rulers}
\label{se:descartes-the-product-of-rulers}

Ren\'e Descartes, in his \emph{La G\'eom\'etrie} of 1637,
introduced the Cartesian coordinate system. A plane is described by
two perpendicular rulers: the $x$-axis and the $y$-axis. Each point
in the plane has two coordinates: its distance along the $x$-axis
and its distance along the $y$-axis. The pair $(x, y)$ identifies
the point.

The construction is the product of two one-dimensional rulers. Each
ruler is a metric record on its own; the plane is the Cartesian
product of the two records. The point $(x, y)$ has identity in
the product: the $x$-coordinate is recovered by the projection
$\pi_x$, the $y$-coordinate by $\pi_y$.

The distance between two points $(x_1, y_1)$ and $(x_2, y_2)$ in
the Cartesian plane is given by the Pythagorean formula:
\[
d = \sqrt{(x_2 - x_1)^2 + (y_2 - y_1)^2}.
\]
This formula is forced by two requirements: (1) the two rulers
are independent, so the distance along $x$ and the distance along
$y$ are admissible separately; (2) the distance is invariant under
rotation of the coordinate system (it does not matter which
direction is ``up'' or ``right'').

The chapter's second claim is that the Cartesian construction is
the canonical reference structure for two-dimensional measurement.
The two perpendicular rulers provide the basis; their product
provides the plane; the Pythagorean formula provides the
distance. The construction extends to higher dimensions: $n$
perpendicular rulers produce an $n$-dimensional space; the
Pythagorean formula generalizes to $d = \sqrt{\sum (x_2^i -
x_1^i)^2}$.

The construction is purely algebraic: a product of one-dimensional
records, with a distance formula determined by the requirement of
invariance under rotation. No geometry is assumed beyond the
metric of the rulers; the plane's geometry is derived from the
metric.

The construction also reveals a subtlety. The Pythagorean formula
depends on the rulers being perpendicular (orthogonal). For
oblique rulers (not at right angles), the distance formula is
different. The choice of perpendicular rulers is the calibration
of the coordinate system: it determines what distance means in
this coordinate system.

For curved surfaces, the Cartesian construction fails: there is no
single pair of perpendicular rulers that covers the whole surface
(the earth's surface, for example, admits no global Cartesian
coordinates). The metric must be defined locally; the metric
tensor (next sections) is the structural object that does this.

\section{The Calibration Certificate}
\label{se:the-calibration-certificate}

Two rulers made of different materials behave differently under
changes of temperature. A steel ruler expands at $11.7 \times
10^{-6}$ per degree Celsius; a brass ruler expands at $19 \times
10^{-6}$ per degree Celsius. At a reference temperature of
$20^\circ$C, the two rulers can be calibrated to agree: both
measure $1$ meter as $1$ meter. At $40^\circ$C, the steel ruler
has grown by $0.000234$ meter; the brass ruler has grown by
$0.00038$ meter; they no longer agree.

To compare measurements made by the two rulers, the calibration
must be named: the temperature at which the calibration holds.
Without the calibration certificate, the rulers' readings cannot
be compared. With it, the readings can be adjusted to account for
the temperature dependence of each material.

The chapter's third claim is that calibration is a structural
prerequisite for comparison across observers. Each observer's
instrument has its own behavior under environmental variations;
the comparison requires naming the calibration conditions. The
calibration is part of the record; the record is admissible
only when the calibration is stated.

In Lean, calibration is the \texttt{EKG} typeclass in
\texttt{LeanCalibration.lean}:

\begin{leanexcerpt}{EKG}
structure EKG where
  private reference : Reading
  embigger? : Prop $\to$ Prop $\to$ Prop
\end{leanexcerpt}

The \texttt{Reading} type is the hidden reference scale (private).
The \texttt{reference} field carries an instance of that scale at
the calibration's baseline. The \texttt{embigger?} method takes
two propositions and returns a proposition asserting that the
first dominates the second under the calibration's bookkeeping.
The calibration's scale is not directly exposed; the comparisons
through the calibration are. The discipline is the chapter's: the
reference is hidden; the comparisons are admissible.

The \texttt{LOGICAL} typeclass requires that a program carry its
\texttt{EKG}:

\begin{leanexcerpt}{LOGICAL}
class LOGICAL
    (Value : Type)
    (Carrier : CarrierProcess Value)
    [DISTINGUISHABLE Value Carrier] ... [UNIVERSAL Value Carrier]
  where
  feelings : HeartbeatProcess Value Carrier
  ekg      : Calibration.EKG
  logical? : YarnTheory $\to$ YarnTheory $\to$ Prop := fun a b =>
    YarnTheory.le a b
\end{leanexcerpt}

A program that asserts comparison results without providing a
\texttt{LOGICAL} instance is not admitted by the compiler. The
calibration certificate is required at every comparison.

The chapter's claim is that this is the mathematical structure of
the calibration discipline. Two observers comparing their counts
must share a reference structure (the \texttt{EKG} or its
analog). The reference structure is hidden behind a public
comparison interface. The comparison's admissibility depends on
the reference structure being correctly calibrated.

For the steel and brass rulers, the calibration certificate names
the temperature. For the SI unit of length (the meter), the
calibration is the definition: the distance traveled by light in
vacuum in $1/299792458$ second. For the chapter's mathematical
records, the calibration is the typeclass \texttt{LOGICAL}
asserting the \texttt{EKG}.

\section{The Metric Tensor}
\label{se:the-metric-tensor}

For a curved surface like the earth, no single global Cartesian
coordinate system exists. The distance between two nearby points
must be computed locally, using the surface's local metric. The
metric tensor $g_{ij}(x)$ at a point $x$ is the structural object
that specifies the local metric.

In local coordinates, the distance between two points with
coordinates $x$ and $x + dx$ is:
\[
ds^2 = g_{ij}(x) \, dx^i \, dx^j = g_{11}(x) (dx^1)^2 + 2 g_{12}(x)
\, dx^1 \, dx^2 + g_{22}(x) (dx^2)^2 + \ldots,
\]
where the sum is over indices $i, j$ ranging over the surface's
dimension. The metric tensor's components $g_{ij}$ vary with
position; the metric is not globally constant for curved surfaces.

For the earth's surface, in spherical coordinates $(\theta,
\phi)$ where $\theta$ is the polar angle and $\phi$ is the
azimuthal angle, the metric is:
\[
ds^2 = R^2 \, d\theta^2 + R^2 \sin^2\theta \, d\phi^2,
\]
where $R$ is the earth's radius. The metric depends on $\theta$:
the longitudinal distance $R \sin\theta \, d\phi$ varies with
latitude. At the equator ($\theta = \pi/2$), the longitudinal
distance is $R \, d\phi$ (the equator's length per unit angle).
At the pole ($\theta = 0$ or $\pi$), the longitudinal distance
is $0$ (a degree of longitude at the pole is zero length).

The chapter's fourth claim is that the metric tensor is the local
reference structure on a curved manifold. The tensor varies with
position; the variation captures the manifold's geometry.
Different observers in different regions use different local
metrics; the comparison across observers requires the metric
tensor's transformation properties under coordinate changes.

In Lean, the metric tensor is part of \texttt{Mathlib}'s
differential-geometry library. The type
\texttt{Manifold.Metric} carries the tensor field; the metric
is computed locally; the geodesic equation and the Levi-Civita
connection are derived from it.

The \texttt{UniverseTensor} type in \texttt{Episode10.lean} is
the chapter's algebraic analog. It carries the local metric at
each point of an admissible record:

\begin{leanexcerpt}{UniverseTensor}
structure UniverseTensor
    (Value : Type)
    (Carrier : CarrierProcess Value)
    [DISTINGUISHABLE Value Carrier] ... [WITNESSED Value Carrier]
  where
  frame_of_reference : ReligiousProcess Value Carrier
  reality            : Truth
  observe?           : Truth $\to$ Truth := ...
\end{leanexcerpt}

The \texttt{UniverseTensor} carries the full cascade through
\texttt{WITNESSED}, a frame-of-reference religious process, a
current \texttt{Truth}, and the function \texttt{observe?} that
advances a \texttt{Truth} through the observer's tensor. The
metric interpretation is metaphorical here: each layer of the
class cascade plays the role of a metric component, and the
\texttt{observe?} function plays the role of the geodesic
projection.

The chapter's claim is that the metric tensor is the most general
form of the reference structure. The Cartesian plane has a
constant metric (the identity matrix). The earth's surface has a
position-dependent metric. General relativity's spacetime has a
position-dependent and signature-indefinite metric. Each is the
chapter's reference structure for its corresponding manifold.

The metric tensor determines distances, angles, areas, and
volumes on the manifold. It determines geodesics (locally
shortest paths). It determines parallel transport (Chapter 7's
connection is derived from the metric). The metric is the
foundational object of Riemannian geometry; it is the chapter's
reference structure.

\section{Truth as Carried Structure}
\label{se:truth-as-carried-structure}

``$2 + 2 = 4$.'' The statement is true. But true relative to what?
Under the Peano axioms for natural numbers, the statement is a
theorem: it can be proved from the axioms by the standard
recursive definition of addition. Under a different axiom system
(say, modular arithmetic in $\mathbb{Z}/3\mathbb{Z}$), the carried
value of the numeral $4$ changes: $4$ now denotes the residue $1$,
and the same arithmetic is written $2 + 2 = 1$. The statement
$2 + 2 = 4$ remains formally true if one interprets $4$ as its
residue class; the disagreement is not in the arithmetic but in
what the symbols carry.

Truth is therefore not absolute; it is relative to the axiom
system. The truth of ``$2 + 2 = 4$'' is a carried certificate:
the proof of the statement from the Peano axioms is what makes
it true. Without the axioms (without the proof carrying the
certificate), the statement is not yet admissible.

The chapter's fifth claim is that mathematical truth is carried
structure. A theorem is true not in some abstract universal sense,
but relative to its proof and the axioms it depends on. The proof
is the certificate; the certificate carries the truth.

In Lean, this is the philosophy of type theory. A proposition $p$
is a type; a proof of $p$ is a term of that type. The
proposition's truth is the existence of the term. Without the
term, the proposition is unproved (and therefore not admissibly
true).

\begin{leanexcerpt}{Truth}
inductive Truth
  | logic : Prop $\to$ Truth
  | fact  : Gospel $\to$ Prop $\to$ Truth $\to$ Truth
\end{leanexcerpt}

The \texttt{Truth} inductive in \texttt{Episode10.lean} carries
the proof history. The \texttt{logic} constructor sits at the
base: a single \texttt{Prop} that grounds the chain. The
\texttt{fact} constructor extends a prior \texttt{Truth} by
attaching a \texttt{Gospel} (the witness from an earlier
chapter), the \texttt{Prop} being added, and the prior truth
stack. The certificate is the entire stack; the truth is portable
exactly as far as the stack's witnesses are portable.

The chapter's claim is that truth is structural, not metaphysical.
Removing the axioms removes the certificate. The certificate is
as portable as the proof that witnesses it. The proof's
portability depends on the axioms being accepted in the destination
context.

This is the chapter's most philosophical claim. It says: truth
is not a cosmic fact but a derived structural property of the
proof system. The chapter does not deny that some things are
true and others false; it asserts that the truth of a specific
statement is relative to the proof system. ``$2 + 2 = 4$'' is
true in the proof system of Peano arithmetic; the truth is
carried by the proof from the Peano axioms.

For the chapter's transport theme, this matters because
transport across different axiom systems is non-trivial. A
theorem proved in one system may not be provable in another.
The transport requires the source system's axioms to be
admissible in the destination system. The transport's
admissibility is checked at the meta-level: do the axioms
agree? If yes, transport is possible; if no, the truth does
not transport.

The chapter closes here. The metric is built from counting.
Descartes's product of rulers gives the Cartesian plane. The
calibration certificate names the reference conditions. The
metric tensor is the local reference structure on curved
manifolds. Truth is carried structure: the proof is the
certificate.

\begin{bridge}

The chapter's question was what reference structure lets distinct
observers compare their counts as one physical record.

The answer is the metric and the calibration certificate. The
metric is built from counting (Hamming distance is the simplest
example). Descartes's product of rulers gives the Cartesian
plane and its Pythagorean distance. The calibration certificate
names the conditions under which two observers' counts can be
compared. The metric tensor is the local reference structure on
curved manifolds. Truth is carried structure: the proof
certificate makes the truth admissible.

What remains is curvature. The metric tensor varies across a
manifold; the variation is the manifold's curvature. The
transport (Chapter 7) accumulates holonomy proportional to the
enclosed curvature. Chapter 9 takes up the question: what
records the failure of transport to commute?

\end{bridge}

\begin{coda}{The Currency Exchange Desk}

The currency exchange desk is in the main terminal of an international
airport. Behind the counter, the clerk has a small computer that
displays the current exchange rates. Above the counter, a board lists
the rates: dollar to euro, euro to yen, dollar to yen, pound to
dollar, and so on. The rates change throughout the day.

A traveler approaches with a hundred-dollar bill. He wants to
convert it to euros. The clerk consults the board. Today's
rate is 0.92 euros per dollar. The clerk calculates: 100 dollars
times 0.92 equals 92 euros. The clerk hands the traveler 92
euros (minus a small commission fee). The transaction is
complete.

The exchange rate is the chapter's reference structure for
currencies. Each currency is its own unit of count; converting
between currencies requires a named reference rate; the rate is
the calibration certificate for the conversion.

Two travelers, one with dollars and one with yen, want to
compare their wealth. The dollars and yen are not directly
comparable: they are counts in different currencies. To
compare, they must convert to a common currency. The common
currency is the chapter's reference structure: the metric in
which both wealths can be expressed.

The exchange rates form a small mathematical structure. Let
$r_{ij}$ denote the exchange rate from currency $i$ to currency
$j$ (the amount of currency $j$ one unit of currency $i$ buys).
For the rates to be admissible, certain consistency conditions
must hold:
\begin{itemize}
\item Identity: $r_{ii} = 1$ for all currencies $i$.
\item Inverse: $r_{ij} \cdot r_{ji} = 1$ (within a small
commission spread).
\item Triangle: $r_{ij} \cdot r_{jk} = r_{ik}$ (the rate from $i$
to $k$ via $j$ equals the direct rate, within a spread).
\end{itemize}

These conditions are the calibration. Without them, the rate
matrix would be inconsistent: an arbitrage opportunity would
exist (buy currency cheaply in one path, sell expensively in
another, profit). The market enforces consistency through
arbitrage: traders identify any inconsistency and trade against
it until the rates align.

The chapter's claim is that this is the structure of any
reference system. The exchange rate matrix is the local metric.
The currency at each city is the local ruler. The consistency
conditions (identity, inverse, triangle) are the calibration
requirements. The arbitrage is the verification that the
calibration holds.

For physical units, the analogous structure is the SI system.
The base units (meter, kilogram, second, ampere, kelvin, mole,
candela) are the reference rulers. The conversion between units
of the same dimension (meter to inch, kilogram to pound) is the
analog of currency exchange. The conversions are calibrated:
1 inch = 0.0254 meter exactly, by definition. The calibration
is internationally agreed, enforced by metrology institutes.

For curved spaces, the analog is the metric tensor at each
point. The tensor at each point is the local ruler; the
relations between tensors at different points are the local
calibration; the global consistency is the manifold's
smoothness.

The traveler walks away with his 92 euros. He has performed the
chapter's transport: a record (his hundred dollars) has been
transported to a different currency (92 euros) through the
calibrated exchange rate. The transport is admissible because
the rate is the calibration; the comparison between currencies
is meaningful because the rate translates one currency's count
into the other's.

The clerk closes out the desk at the end of her shift. The
day's transactions are logged: the conversion volumes, the
realized exchange rates, the commission income. The log is
the chapter's record of the day's reference activity. The
discipline of the currency exchange is the chapter's
discipline in financial form: every comparison requires a
named reference; the reference must be consistent; the
consistency is checked at the level of the rate matrix.

\end{coda}
